\documentclass[10pt,a4paper]{amsart}
\usepackage[margin=19mm]{geometry}
\usepackage{amsmath,amssymb,mathtools}
\usepackage[scr=rsfs,cal=euler]{mathalfa}
\usepackage{xfrac}
\usepackage[unicode,hidelinks,bookmarks=false]{hyperref}
\newcommand{\ie}{i.e.\ }
\newcommand{\cf}{cf.\ }
\newcommand{\resp}{respectively\ }
\newcommand{\Subsection}{\subsection}
\providecommand{\nobreakdash}{\nobreak\hbox{-}}
\setlength{\emergencystretch}{2em}
\multlinegap0pt
\usepackage[shortlabels]{enumitem}
\usepackage{amssymb}
\newtheorem{lemma}{Lemma}
\title{A proof of purely singular splitting conjecture}
\author{Ka Hin Leung, Tao Zhang}
\date{Source: arXiv:2605.09871v1 (CC BY 4.0)}
\begin{document}
\maketitle

\noindent\emph{Source and licence.} A proof of purely singular splitting conjecture. Version arXiv:2605.09871v1 (CC BY 4.0), distributed by arXiv under the Creative Commons Attribution 4.0 International licence, http://creativecommons.org/licenses/by/4.0/, as recorded in the arXiv licence metadata for this exact version and shown on the abstract page https://arxiv.org/abs/2605.09871v1. Full licence notice: \texttt{licenses/arxiv-2605.09871-CC-BY-4.0.txt}.\par\smallskip

\noindent\textit{Editorial scope.} Counted unit: the article's lemma lemma--1 (source0.tex lines 211--217). The article's complete original proof of it is delivered with it (source0.tex lines 219--244, one \texttt{proof} environment of 26 lines). No mathematics is written, repaired or re-derived: the statement and the proof are the article's own lines, in the article's own notation, and no retained line is substituted or re-flowed.  No further source-local context is needed: every cross-reference of every retained line resolves inside the two delivered spans.  Nothing was omitted from any retained span, so the package carries no omission record.  No retained line contains a citation, so the wrapper carries no bibliography and no bibliography entry.  No retained line contains a figure environment, an image-inclusion command or drawing code, so no image is delivered and the package declares no asset dependency.  Editorial formatting only, in addition: an AMS \texttt{amsart} preamble that restates the article's own declarations and macros used by the retained lines, namely the environments \texttt{lemma} on separate counters, together with the standard packages those lines need.  The retained environments print in this document as Lemma 1 in order of appearance, whereas the article prints the same items with the numbering of its own full document.  The complete native source of the article as distributed in the arXiv e-print for this exact version is bundled unchanged as \texttt{sources/200267/source0.tex}.
\begin{lemma}\label{lemma--1}
	\begin{itemize}
		\item[(a)] \(TW_i \cap TW_j = \emptyset\) for all distinct \(1 \le i,j \le r\).
		\item[(b)] $\bigcup_{i=1}^{r} TW_i \subseteq \mathbb{Z}_{p^\alpha m}^{*},$
		and  $r|TW_i|\le|\mathbb{Z}_{p^\alpha m}^{*}|.$
	\end{itemize}
\end{lemma}

\begin{proof}
	(a) Suppose, to the contrary, that there exist \(t_1,t_2\in T\) and \(x_1,x_2\in M'\) such that
	\[
	t_1(x_1+s_i)=t_2(x_2+s_j).
	\]
	Multiplying both sides by \(p^\beta m'\), we obtain
	\[
	p^\beta m' t_1(x_1+s_i)
	=
	p^\beta m' t_2(x_2+s_j).
	\]
	Since \(p^\beta m' x_1=p^\beta m' x_2=0\) in \(G\), it follows that
	\[
	p^\beta m' t_1 s_i
	=
	p^\beta m' t_2 s_j.
	\]
	Observe that $1\le p^\beta m' t_1,\ p^\beta m' t_2 \le k.$
	Hence the above equality contradicts the uniqueness of representations in the splitting. Therefore, $TW_i \cap TW_j=\emptyset$
	whenever \(i\ne j\).
	
	(b) By the definition of \(W_i\), we have $W_i\subseteq \mathbb{Z}_{p^\alpha m}^{*}.$
	Since both \(T\) and \(W_i\) are subsets of \(\mathbb{Z}_{p^\alpha m}^{*}\), it follows that $\bigcup_{i=1}^{r} TW_i
	\subseteq\mathbb{Z}_{p^\alpha m}^{*}.$
	The inequality $r|TW_i|\le|\mathbb{Z}_{p^\alpha m}^{*}|$ now follows immediately from part~(a).
\end{proof}

\end{document}
